% 04-discussoes.tex | 4 Discussions (nome exigido pela Production)
% RESTAURADO DO ORIGINAL em 06/11/2026: secao "4. Discussion of Results" do Paper Training - 15 Nov 2025
% (reacao, aprendizagem, comportamento, resultados, sintese), quase sem cortes.
% Da versao JATM (pedidos 1 e 2 do revisor do JATM): implicacoes gerenciais (4.6) e agenda para
% reguladores (4.7), sem os numeros que nao tem origem (custo de implantacao de USD 80 a 110 mil,
% payback, ROI) e com as recomendacoes regulatorias como agenda a testar (parecer de 11/08/2026).
% Citacoes em vermelho = referencias do original nao localizadas.

\section{Discussions}
\label{sec:discussions}

This study operationalized Kirkpatrick's four-level model within a high-stakes aviation \MRO{} environment, illustrating its dual utility as a tool for safety enhancement and for sustainability performance. The findings indicate that a structured training evaluation framework is viable under the stringent constraints of \AS{} and civil aviation regulations and is also a mechanism for operational excellence and resource efficiency.

\subsection{Reaction level: foundations for engagement and efficiency}
\label{subsec:disc-reaction}

The positive reception captured at the reaction level, together with the identification of specific infrastructure and content gaps (Figure~\ref{fig:scores-reaction}), provides the foundational data for continuous improvement. More than a measure of satisfaction, this level acts as an early indicator of systemic resource inefficiencies. Feedback on outdated training facilities or unclear content points to areas where suboptimal resource allocation can lead to knowledge gaps that later appear as physical waste, such as misassembled components that require rework or scrap. Addressing these issues early is a first step in a waste-prevention strategy within the training value chain, beyond its value for employee satisfaction.

\subsection{Learning level: verifying knowledge to prevent future waste}
\label{subsec:disc-learning}

The learning assessment, which identified specific topics requiring clarification (mean 66.90 against a pass mark of 70), works as a quality checkpoint before knowledge gaps turn into errors on the shop floor. The gaps identified are potential failure points rather than academic shortcomings: they can lead to future non-conformances and material loss. By verifying knowledge acquisition, the organization can target its instructional resources more effectively, preventing the downstream economic and environmental costs of reworking or scrapping high-value engine parts because of procedural misunderstandings \parencite{Chen2023}. This turns the learning evaluation from a simple test into a proactive risk-mitigation tool for both safety and resource conservation.

\subsection{Behavior level: translating knowledge into sustainable practices}
\label{subsec:disc-behavior}

The behavioral changes and error reduction observed at this level are the most direct evidence of the training's effect on operational performance. The reduction in errors is a safety milestone and also a sustainability gain: each avoided error represents a conserved resource, a saved component, hours of rework labor prevented and energy consumption averted in test cells. The use of two perspectives, self-reported and supervisory, strengthens this finding and indicates that the training migrated from the classroom to the hangar, encouraging work practices that are safer and less wasteful \parencite{YangPark2024}. \todo{sustentar com os resultados do Nível 3 (médias das 15 duplas), que o original não traz}

\subsection{Results level: quantifying the strategic impact on performance and waste}
\label{subsec:disc-results}

The customer satisfaction ratings show where the preceding levels should act: technical documentation and turnaround time were the weakest categories (rated 3) in the first months of the study year, and repeated surveys can capture the cumulative effect of the four levels. Better quality means fewer external failures, which are among the costliest forms of waste, involving warranty claims, concessions and large resource outlays for off-site repairs. A shorter turnaround time would reflect a smoother, less resource-intensive process. These results-level metrics suggest that effective training contributes to the triple bottom line: economic performance, by reducing the cost of quality; social performance, through improved safety; and environmental performance, by minimizing the material and energy footprint of maintenance operations \parencite{Sousa2025}. \todo{não há medida de satisfação anterior à intervenção (Fig. 8 é de um período só); falar em níveis de satisfação ou acrescentar o dado anterior} Because the economic estimate rests on the same reduction in training failures, it re-expresses that reduction in monetary terms rather than measuring it independently.

\subsection{Synthesis: a framework for integrated compliance and sustainability}
\label{subsec:disc-synthesis}

A strength of this study is that it illustrates how the Kirkpatrick model provides the structure needed to meet and exceed the \AS{} and CAR mandates for objective evidence of competence. The framework turns abstract regulatory requirements into a clear, data-driven workflow for proving training effectiveness. It also raises training from a compliance activity to a strategic function that can protect human life, improve financial performance and reduce environmental impact at the same time. Each level maps onto \AS{} clauses or an \SMS{} component and onto a class of sustainability indicators (Figure~\ref{fig:conceptual}). This combination of compliance, safety and cleaner production offers a replicable blueprint for \MRO{} organizations in an industry increasingly focused on sustainable aviation.

\subsection{Managerial implications}
\label{subsec:managerial}

For \MRO{} managers, the evaluation fits the \LMS{} and \ERP{} architecture most facilities already run, in five steps:

\begin{enumerate}[label=(\arabic*)]
  \item a mandatory Level~1 exit form in the \LMS{};
  \item a Level~2 assessment with a pass mark (70/100 or a local equivalent) that conditions \ERP{} access;
  \item supervisors trained in the Level~3 assessment, with records in each technician's training file;
  \item Level~4 metrics (training failures, regulatory non-conformances, quality costs, customer satisfaction) as a standing item of the monthly quality meeting;
  \item Level~3 and Level~4 data used to revise the next training plan.
\end{enumerate}

The same records support audits. Level~2 and Level~3 records evidence the clause 7.2 requirement to evaluate training effectiveness; Level~1 records and the training plan support clauses 7.1, 7.3 and 7.4; and Level~4 trends support clause~9. For sustainability reporting, they can feed indicators such as avoided material waste, avoided rework hours, avoided expedited logistics (kgCO\textsubscript{2}e) and training hours per direct labor hour.

\subsection{Agenda for regulators and industry bodies}
\label{subsec:policy}

A single site cannot ground mandates, but the case suggests three propositions to test in multi-site pilots:

\begin{enumerate}[label=(\arabic*)]
  \item civil aviation authorities could ask for evidence of four-level training evaluation when accepting \SMS{} manuals, since \ICAO{} Annex~19 leaves the depth of evaluation to the service provider \parencite{Nakamura2024,OConnor2024};
  \item industry bodies such as IATA and SAE could issue joint guidance on training evaluation for \MRO{}, with minimum data sources, metrics by level and a maturity model;
  \item regulators could trial, as proposed here, confidential annual reporting of Level~3 and Level~4 metrics as safety performance indicators, which \SMS{} safety assurance already tracks \parencite{FAA2024}, linking them, where environmental, social and governance reporting applies, to the human-capital sections of sustainability reports.
\end{enumerate}
